\section{Projet 17 — Memory}
\setcounter{UPSTInumeroCahier des charges}{0}

\subsection*{Objectif du jeu}
Retrouver les 8 paires de valeurs identiques cachées dans une grille de 16 cases, en un minimum de tentatives.

\subsection*{Règles du jeu}
\begin{itemize}
  \item La grille compte \textbf{16 cases} ($4\times4$, tableau d'entiers ou de \texttt{char}) contenant \textbf{8 paires} : les valeurs \texttt{1} à \texttt{8}, chacune en double, placées au hasard.
  \item Par défaut toutes les cases sont \textbf{cachées} et affichées \texttt{?}.
  \item À chaque tentative, le joueur indique \textbf{deux positions distinctes} ; leurs valeurs sont révélées.
  \item Si les deux valeurs sont \textbf{identiques}, la paire est trouvée et reste visible définitivement ; sinon les deux cases sont de nouveau cachées après une courte pause (\texttt{sleep()}, voir section Bonus).
  \item Le jeu est \textbf{gagné} lorsque les 8 paires sont trouvées.
\end{itemize}

\textit{Conseil :} on utilise deux tableaux de même taille : l'un contient les \textbf{valeurs}, l'autre l'\textbf{état} de chaque case (trouvée ou non). Le mélange initial repose sur des échanges de cases, comme le tirage sans remise (voir la fiche d'aide « Techniques utiles », section Tirage sans remise).

\subsection*{Cahier des charges}
\begin{UPSTICahierDesCharges}{1}
\begin{itemize}
  \item Remplir la grille avec les valeurs \texttt{1} à \texttt{8} \textbf{en double}.
  \item \textbf{Mélanger} la grille aléatoirement par échanges de cases.
  \item Afficher la grille \textbf{cachée} (\texttt{?}) ; en mode test, afficher aussi la grille réelle.
\end{itemize}
\end{UPSTICahierDesCharges}

\begin{UPSTICahierDesCharges}{2}
\begin{itemize}
  \item Saisir \textbf{deux positions} (vérifier qu'elles sont valides, distinctes et non déjà trouvées).
  \item Afficher \textbf{temporairement} les deux valeurs réelles, puis les \textbf{comparer}.
  \item Mettre à jour le \textbf{tableau d'état} (case trouvée oui/non) et ré-afficher la grille après la pause.
\end{itemize}
\end{UPSTICahierDesCharges}

\begin{UPSTICahierDesCharges}{3}
\begin{itemize}
  \item Compter et afficher le \textbf{nombre de tentatives}.
  \item Détecter la \textbf{victoire} (8 paires sur 8).
  \item Structurer en fonctions (remplissage, mélange, affichage, saisie, comparaison).
\end{itemize}
\end{UPSTICahierDesCharges}

\textit{Amélioration :} mode \textbf{2 joueurs} : le joueur qui trouve une paire rejoue ; le vainqueur est celui qui a trouvé le plus de paires.

